% Figure: velocity field of a wheel rolling without slipping, showing the
% instantaneous centre at the contact point and velocities at key points.
\begin{figure}[htbp]
\centering
\begin{tikzpicture}[
  vel/.style={draw=engblue, -{Stealth}, very thick},
  lbl/.style={font=\small},
]
% ground
\draw[engblack, thick] (-2.2,0) -- (4.2,0);
% wheel of radius 1.4 centred at (1,1.4)
\coordinate (O) at (1,1.4);
\coordinate (P) at (1,0);     % contact (instantaneous centre)
\coordinate (T) at (1,2.8);   % top
\coordinate (F) at (2.4,1.4); % front (right)
\coordinate (B) at (-0.4,1.4);% back (left)
\draw[draw=engblack, thick, fill=enggray!10] (O) circle (1.4);
\foreach \pt in {O,T,F,B} {\fill[engblack] (\pt) circle (1.3pt);}
\fill[engred] (P) circle (1.8pt);
% velocity at centre: v = omega R, to the right
\draw[vel] (O) -- ++(1.6,0) node[above, lbl] {$\vec{v}_{G}=\omega R$};
% velocity at top: 2 v
\draw[vel] (T) -- ++(3.0,0) node[above, lbl] {$\vec{v}_{\text{top}}=2\omega R$};
% velocity at front: directed up-back at 45 (magnitude sqrt2 v), but
% draw vertical-ish for clarity: actually v_F = omega x r, r from P to F.
% From instantaneous centre P=(1,0): F=(2.4,1.4), v perpendicular to PF.
\draw[vel] (F) -- ++(0.95,-0.95) node[right, lbl] {$\sqrt{2}\,\omega R$};
% velocity at back: from P=(1,0) to B=(-0.4,1.4), perpendicular
\draw[vel] (B) -- ++(0.95,0.95) node[left, lbl] {$\sqrt{2}\,\omega R$};
% contact point: zero velocity, mark
\node[engred, font=\small, below] at (P) {$\vec{v}_{P}=\vec{0}$};
\node[engred, font=\scriptsize, below right] at ($(P)+(0.05,-0.45)$) {instantaneous centre};
% omega arc
\draw[draw=engamber, -{Stealth}, thick] (O) ++(140:0.55) arc (140:-30:0.55);
\node[engamber, font=\small] at ($(O)+(0.0,0.78)$) {$\omega$};
\end{tikzpicture}
\caption[Velocity field of a rolling wheel]{Velocity of selected
points on a wheel rolling without slipping. The contact point $P$
(red) is the instantaneous centre of rotation: its velocity is zero
at that instant, which is exactly the no-slip condition. Every other
point moves perpendicular to its line from $P$, with speed
proportional to its distance from $P$. The centre $G$, at distance
$R$ from $P$, moves at $\omega R$; the top, at distance $2R$, moves
at $2\omega R$; the leading and trailing points, at distance
$\sqrt{2}\,R$, move at $\sqrt{2}\,\omega R$ directed perpendicular to
their radius from $P$. Treating the motion as pure rotation about the
instantaneous centre reproduces the translation-plus-rotation result
$\vec{v}_{B}=\vec{v}_{G}+\vec{\omega}\times\vec{r}_{B/G}$ with less
bookkeeping.}
\label{fig:vol03ch04:rolling-velocity}
\end{figure}
